\PassOptionsToPackage{margin=1in}{geometry}
\documentclass[twocolumn]{NobArticle}
\runninghead{Kidney Yogurt}

% === Packages for clean, modern look ===
%\usepackage[letter,margin=1in]{geometry}          % Reasonable margins
%\geometry{letterpaper, margin=1in}
\usepackage{fontspec}                      % For better fonts (use XeLaTeX or LuaLaTeX)
\setmainfont{Times New Roman}              % Or your favorite: Georgia, Calibri, etc.
\usepackage{parskip}                       % Space between paragraphs instead of indent
\usepackage{titlesec}                      % Nicer section titles
\usepackage{xcolor}                        % Colors for emphasis
\usepackage{hyperref}                      % Clickable links
\hypersetup{colorlinks=true, linkcolor=blue, urlcolor=blue}
\usepackage{booktabs}                      % Beautiful tables
\usepackage{enumitem}                      % Custom lists
\usepackage{graphicx}                      % Images (e.g. LTG photos)
\usepackage{caption}                       % Better captions



% === Section formatting (clean & modern) ===
\titleformat{\section}
{\large\bfseries\color{black!70!black}}
{\thesection}{1em}{}
\titleformat{\subsection}
{\normalsize\bfseries\color{black}}
{\thesubsection}{1em}{}

% === Light page background note (optional) ===
\usepackage{background}
\backgroundsetup{
	scale=1,
	angle=0,
	opacity=0.07,
	contents={Draft - Lily the Great Recovery Journal}
}

% === Title & Metadata ===
% Title
\title{Saddle Yogurt: How a Homemade Probiotic Helped My Cat Recover from Saddle Thrombus}
% Authors
\author{Mike Stephenson, aka Lily the Great's Dad\textsuperscript{1}}
% Affiliations
\date{\textsuperscript{\textbf{1}}Solverware, San Diego, California}


\begin{document}
	
	\maketitle
	
	\begin{abstract}
		This is the story of how I used a homemade probiotic yogurt based on \textit{Bacillus Subtilis} to help my cat Lily the Great recover from a serious saddle thrombus event that left her without use of her hind legs. The saddle thrombus blood clot dissolved quickly based on her activity level which improved every day and resulted in a thriving and happy 18.5 year old cat. The active compound that resulted in this remarkable recovery is the enzyme Nattokinase, which is produced by the fermentation of soybeans by the microbe \textit{Bacillus Subtilis}. The yogurt serves as a delivery mechanism for the Nattokinse supplement pills but also results in \textit{B. Subtilis} colonize the small intestine. \textit{B. Subtilis} will digest anything Lily ate in the small intestine and as such, the yogurt turned Lily into a Nattokinase factory right in her small intestine. It is unclear how much Nattokinase gets produced in this manner so it should be considered as \textit{augmenting} the amount ingested via supplements that makes it past her stomach in tact. This is not veterinary advice—always consult your vet before trying any supplement or diet change. N.B. - this is anecdotal data for one cat so results may vary.
	\end{abstract}
	
	\section{Introduction}
	Lily the Great (LTG) is an 18.5 year old calico cat born in May 2007 in El Centro, California. She relocated to sunny San Diego, in June of 2007 and choose two humans, Linda and Mike, who lived just down the street. She has been a mostly healthy and happy cat bring much pleasure to her humans.
	
	\section{Health History}
	Lily was diagnosed with chronic kidney disease (CKD) back in 2021. We switched her diet to eliminate kibble, which was using to augment her diet for times when we were away for several hours. It was much easier to put out kibble and know she would not go hungry. After studying CKD and getting some counsel, we switched her diet to be primarily raw food but 100\% wet food.
	
	\section{Suspected Saddle Thrombus Events},
	Lily's first event that we believe is saddle thrombus related occurred on November 1, 2024. It left Lily without full control of her hind legs. I was out for the evening at a scheduled event until 9:30. Linda was home. I received a text message from Linda at 9:02 p.m. that read "Please come home ASAP... I'm afraid Lily is having problems." When I saw the message, I phoned immediately at 9:36 p.m. Lily's event had already past, I was home by 10:20 p.m. and all seemed normal.
	
	The second event occurred about six months later in April 2025. It was similar to the first but more persistent and resulted in a decline in Lily's quality of life. She was taken to the veterinarian and nothing was brought up about a possible saddle thrombus event. Because of Lily's kidney disease progression, the vet recommended subcutaneous fluids every other day, which we started doing that. LTG decided at some point after this event to take up residence in the garage rather than the house, which was a marked change in behavior. She hopped up on chairs, couches, less frequently. While she had often slept in the bedroom, the jump to get on the bed was simply too high so she stopped that practice.
	
	The third and most severe event occurred on December 10, 2025. The event occurred sometime between 10:05 a.m. when I left the house and 10:45 when Linda got home. Lily was in some distress and pain, evident by a painful sounding yowl. Lily could not walk or stand up. Lily was rushed to the vets and given an exam and we were debriefed. Linda was there in person and she got me on the phone. The vet suggested tests to figure out what occurred and what to do about it. At this point, all we knew was Lily lost the use of both hind legs and was \textit{effectively} paralyzed. With Lily's kidney disease and the vet already counseling us on euthanasia, and no definitive diagnosis without an appreciably financial outlay, we decided to take her home and reassess. We basically had little data to go on. 
	
	When I put the blood labs into Grok AI, which I had been using to consult with regarding Lily's kidney disease, three of her four kidney markers actually improved. I worked with Grok to develop a so-called "kidney yogurt." We both had noticed up until the event that Lily had a marked resurgence which was 100\% time-coincident with the introduction of the kidney yogurt.
	
	
	% === Your content here ===
	\section{Saddle Thrombus in Cats}	
	Saddle thrombus is 
	
	
	\section{The Microbiome and Super Gut Yogurts}
	The microbiome is the collection (or community) of all microorganisms—including bacteria, archaea, fungi, viruses, and sometimes other tiny eukaryotes—along with their genes and the environmental conditions of the habitat they occupy in a specific location, usually the human body. The human body hosts as many microbe cells as it does human cells and microbial genes outnumber human genes by as much as a factor of 100. 
	
	The microbiome plays huge roles in digestion, immune system training, vitamin production, protection against pathogens, metabolism, and even brain function via the gut-brain axis. When it's out of balance (dysbiosis), it can contribute to many diseases.
	
	Cats also have a microbiome.
	
	\section{The Gut-Kidney Axis}
	Chronic kidney disease (CKD) is very common in older cats (affecting up to ~30-50\% over age 10-15), and the gut-kidney axis is an emerging area of interest. Gut dysbiosis (imbalanced microbiota) can contribute to uremic toxin buildup (e.g., from gut bacteria metabolizing proteins into harmful compounds like urea/ammonia/indoxyl sulfate), worsening kidney burden and inflammation.
		
	\input{saddle-yogurt}
		
	\section{The Science Behind Kidney Yogurt}
	B. coagulans is a spore-forming probiotic...
	%\section{How to Make Kidney Yogurt}


The book Super Gut, by reknowned cardiologist Dr. William Davis describes various homemade yogurt concoctions are made by a long fermentation process, usually in a dairy substance such as organic half-and-half. A distinguishing point of this process is the long fermentation time on the order of 36 hours. The fermentation time is designed to produce high colony forming units (CFUs) of whatever probiotic is being fermented. These microbes reproduce asexually so they double according to a time interval, often times as short as 30 minutes but can be 3 or more hours. The double time tends to be a function of temperature and the fermentation environment, including the medium (e.g., half-and-half).

The microbe that has perhaps gotten the most attention is L. Reuteri as it appears to be a keystone microbe that has mostly been lost in modern society due to overuse of antibiotics. L. Reuteri has a double time of three hours which equates to 12 doublings in a 36-hour fermentation cycle.

Subtilis doubles every 20-25 minutes at 37°C. 

\begin{table}[ht]
	\centering
	\begin{tabular}{ll}
		\toprule
		\textbf{Ingredient} & \textbf{Amount} \\
		\midrule
		Organic ultra-pasteurized half-and-half & 1 quart \\
		B. coagulans GBI-30,6086 capsules & 1--2 \\
		Inulin powder (prebiotic) & 2 tbsp \\
		\bottomrule
	\end{tabular}
	\caption{Ingredients for 1 quart batch}
	\label{tab:recipe-ingredients}
\end{table}

Fermentation steps:
\begin{enumerate}[leftmargin=*]
	\item Mix slurry...
	\item Heat and cool...
	\item Ferment at 106°F for 36 hours...
\end{enumerate}
	
	The book Super Gut, by reknowned cardiologist Dr. William Davis describes various homemade yogurt concoctions are made by a long fermentation process, usually in a dairy substance such as organic half-and-half. A distinguishing point of this process is the long fermentation time on the order of 36 hours. The fermentation time is designed to produce high colony forming units (CFUs) of whatever probiotic is being fermented. Such probiotics reproduce asexually so double based on a time factor, which is also a function of temperature and the fermentation environment.
	
	The microbe that has perhaps gotten the most attention is L. Reuteri as it appears to be a keystone microbe that has mostly been lost in modern society due to overuse of antibiotics. L. Reuteri has a double time of three hours which equates to 12 doublings in a 36-hour fermentation cycle.
	 
	Subtilis doubles every 20-25 minutes at 37°C. 
	
	\begin{table}[ht]
		\centering
		\begin{tabular}{ll}
			\toprule
			\textbf{Ingredient} & \textbf{Amount} \\
			\midrule
			Organic ultra-pasteurized half-and-half & 1 quart \\
			B. coagulans GBI-30,6086 capsules & 1--2 \\
			Inulin powder (prebiotic) & 2 tbsp \\
			\bottomrule
		\end{tabular}
		\caption{Ingredients for 1 quart batch}
		\label{tab:recipe-ingredients}
	\end{table}
	
	Fermentation steps:
	\begin{enumerate}[leftmargin=*]
		\item Mix slurry...
		\item Heat and cool...
		\item Ferment at 106°F for 36 hours...
	\end{enumerate}
	
	\section{Lily's Journey \& Lab Results}
	Before yogurt (April 2025): Creatinine 9.8, BUN >130...  
	After 7 weeks: Creatinine 6.7...
	
	\section{Conclusion}
	If your cat has CKD, discuss probiotics with your vet...
	
	\printbibliography
	
\end{document}